\documentclass[11pt]{article}
\usepackage{icmlko}
\begin{document}

\title{축소 모형에서 잰 것을 챔피언에 넣었다 --- 노트 441$\sim$443 을
되돌린다}
\author{월드모델 실험실}
\date{노트 463 · 2026년 8월 2일}
\maketitle

\begin{abstract}
노트 462 의 진단이 \texttt{OOC\_DROP} 을 꺼야 해서 껐는데, 그 팔의 KR 이
챔피언보다 \emph{높게} 나왔다. 확인해 보니 우연이 아니었다.

\begin{center}
\begin{tabular}{lrrr}
\hline
짝 & 가림 켬 & 가림 끔 & 차(켬$-$끔) 95\% \\
\hline
KR 만화 & $+0.5536$ & $\mathbf{+0.6119}$ & $-0.0583$ $[-0.0611,-0.0555]$ \\
비게임 앱 & $\mathbf{+0.2646}$ & $+0.2232$ & $+0.0414$ $[+0.0296,+0.0532]$ \\
\hline
\end{tabular}
\end{center}

\textbf{두 짝이 갈렸다.} 노트 441 이 채택 근거로 삼은 조건 ①(집 밖 짝
둘이 같은 부호)이 지금 챔피언에서 안 선다.

챔피언 조각 탓인가 하고 하나씩 껐다 --- 깊이 $6$ 으로 되돌리고 계절과
특기 칸을 다 꺼도 KR 의 가림 이득은 $-0.0387$ 로 \textbf{여전히 음수}다.

\textbf{남은 차이는 내 실수였다.} 노트 $441 \sim 443$ 의 측정은 전부
\textbf{AX5(+gen/grp)만 학습에 넣은 축소 모형}에서 했다. 구현은
\textbf{36축 챔피언}에 넣었고, \textbf{채택 뒤 챔피언에서 이득이 남아
있는지 다시 안 쟀다.} 그때 붙인 가드 셋은 집안 예보 불변 · 가지 열림 ·
입력 불변만 봤다 --- \emph{작동하는지}는 봤고 \emph{이로운지}는 안 봤다.
\end{abstract}

\begin{figure}[t]
\centering
\includegraphics[width=\textwidth]{figs/f1.pdf}
\caption{같은 세 팔을 축소 모형(가)과 챔피언(나)에서 잰 것. 초록 삼각형이
그 짝의 최선. \textbf{KR 의 순서가 완전히 뒤집힌다} --- 축소 모형에서는
뒤집기가 최선이고 그대로가 최악인데, 챔피언에서는 정반대다.}
\end{figure}

\section{노트 443 은 반증됐다}

축이 다섯일 때 노트 443 은 \emph{``라벨 다섯이면 버린 축을 되찾고 마흔이면
오라클''}이라 적었다. 36축에서 다시 쟀다(씨앗 여덟 짝지어).

\begin{center}
\begin{tabular}{lrrr}
\hline
짝 & 그대로 & 가림 & 뒤집기 \\
\hline
KR 만화 & $\mathbf{+0.6115}$ & $+0.5531$ & $+0.3833$ \\
비게임 앱 & $+0.2213$ & $\mathbf{+0.2648}$ & $+0.2474$ \\
\hline
\end{tabular}
\end{center}

\textbf{뒤집기가 두 짝 어디에서도 최선이 아니다.} KR 에서는 가림보다도
$-0.1698$ $[-0.1781,-0.1615]$ 나쁘다. 그런데 \texttt{calibrate\_ooc} 는
짝 안 부호가 음수면 뒤집으므로, KR($-0.320$)에 부르면
\textbf{$+0.6115$ 를 $+0.3833$ 으로 떨어뜨린다.} 안 부르고 있었을 뿐,
\emph{부르면 해로운 메서드}를 두 사이클째 들고 있었다.

\textbf{읽는 법.} 축이 다섯뿐이면 \texttt{entry\_friction} 이 입력의 큰
몫이라 그 축의 부호가 사실상 전부였다. 36축에서는 모델이 그 축을 다른
것들과 \emph{함께} 쓰므로, 값을 뒤집으면 배운 조합이 깨진다. \textbf{부호를
고치는 처방은 축이 적을 때만 옳다.}

\section{무엇이 남고 무엇이 죽나}

\begin{center}
\begin{tabular}{ll}
\hline
노트 441 --- \texttt{entry\_friction} 만 부호가 안 맞는다 & \textbf{산다}(자료 사실) \\
노트 441 --- 안 배운 수집에 가린다 & \textbf{미결}(두 짝이 갈림) \\
노트 442 --- 계산되는 표지가 원핫보다 나쁘다 & \textbf{미결}(축소 모형에서 잼) \\
노트 443 --- 라벨 몇 개로 뒤집는다 & \textbf{반증} \\
\hline
\end{tabular}
\end{center}

\texttt{OOC\_DROP} 은 \textbf{유지}한다 --- 켬도 끔도 조건 ①을 못 채우고,
판은 이 변경을 원리상 못 보며(집안 예보가 정확히 같다), 노트 444 의
천장 비로 정규화해도 $60.1\%$ 대 $60.9\%$ 로 무차별이다. \textbf{유지하되
근거를 ``증명''에서 ``미결''로 내린다.}

\texttt{calibrate\_ooc} 는 코드에 남기되 \textbf{뒤집기를 기본으로
껐다}(\texttt{allow\_flip=False}). 부호는 계속 재서 적어 둔다 --- 언젠가
축이 적은 수집을 만나면 그때 켤 근거가 생길 수 있다.

\section*{규약 --- 채택한 뒤 챔피언에서 다시 잰다}

\begin{quote}
\textbf{조항} --- 축소 모형 · 부분 설정에서 잰 효과는 \textbf{구현 뒤
챔피언에서 다시 재기 전까지 ``미결''로 적는다.} 채택 때 붙이는 가드는
\emph{작동하는지}를 보는 것이고, \emph{이로운지}는 따로 재야 한다.
\end{quote}

이 사이클에서 규약이 세 번 일했다 --- 노트 451 이 내 문턱을 고쳤고,
노트 458 이 채택 직전을 막았고, 노트 460 이 통과한 거래를 환율로 막았다.
\textbf{이번엔 규약이 없어서 두 사이클을 잃었다.} 없던 자리에 하나 더
박는다.

\section*{남은 것}

\begin{itemize}
\item \textbf{노트 442 를 36축에서 다시 재기} --- 계산되는 표지가 정말
원핫보다 나쁜가. 미결로 남았다.
\item \textbf{다른 축소 측정이 또 있나} --- 이 사이클에서 짝 실험은
대부분 AX5 로 돌렸다. 어느 결론이 축 수에 기대는지 훑어야 한다.
\item \textbf{영샷 기본값을 정할 셋째 짝} --- 두 짝이 갈리는 한
\texttt{OOC\_DROP} 은 계속 미결이다.
\end{itemize}

\end{document}
